\documentclass[12pt,a4paper]{article}

\usepackage[T2A]{fontenc}
\usepackage[utf8]{inputenc}
\usepackage[russian,english]{babel}
\usepackage{geometry}
\usepackage{graphicx}
\usepackage{booktabs}
\usepackage{longtable}
\usepackage{array}
\usepackage{multirow}
\usepackage{amsmath}
\usepackage{amssymb}
\usepackage{siunitx}
\usepackage{xcolor}
\usepackage{hyperref}
\usepackage{caption}
\usepackage{subcaption}

\geometry{margin=2.2cm}
\graphicspath{{../reports/figures/article_real/}}
\hypersetup{
    colorlinks=true,
    linkcolor=blue!50!black,
    citecolor=blue!50!black,
    urlcolor=blue!50!black
}

\sisetup{
    output-decimal-marker={,},
    group-separator={\,},
    group-minimum-digits=4
}

\title{Блок 1. Полные результаты: достоверность цифрового двойника и разрыв между прогнозной точностью и полезностью для RL}
\author{HVAC DRL MORL Project}
\date{\today}

\begin{document}
\selectlanguage{russian}
\maketitle

\begin{abstract}
Блок 1 формирует эмпирическое основание для центрального научного утверждения проекта: прогнозная точность суррогатной модели и её полезность как среды обучения reinforcement learning являются разными свойствами. В работе сравниваются два суррогата тесткейса BOPTEST \texttt{bestest\_air}: компактный control-oriented суррогат v3 и физически информированный суррогат v3.5 с обратной калибровкой Stage A/B/C. Калиброванный v3.5 улучшает 24-часовой rollout RMSE температуры с \SI{1.466}{\celsius} до \SI{0.644}{\celsius}, а параметр тепловой ёмкости зоны \(C_{\mathrm{zon}}\) идентифицируется как \(4.413\times 10^5\) Дж/К, что всего на \SI{5.1}{\percent} выше физически мотивированного prior. Однако при использовании v3.5 как прямой среды обучения PPO live BOPTEST результат ухудшается: RMSE замкнутого контура достигает около \SI{4.32}{\celsius}, тогда как v3 даёт около \SI{0.89}{\celsius}. Гибридная схема разрешает этот парадокс: v3 используется как гладкая rollout-динамика, а калиброванный v3.5 --- как замороженный физический регуляризатор. Итоговый вывод Блока 1: более точный цифровой двойник не обязан быть лучшей RL-средой; физическая модель наиболее полезна как teacher/regularizer, а не как прямой rollout backend.
\end{abstract}

\tableofcontents
\newpage

\section*{Executive Summary}
\addcontentsline{toc}{section}{Executive Summary}

Блок 1 проверяет три вопроса.

\begin{enumerate}
    \item Можно ли построить быстрый control-oriented суррогат для обучения PPO?
    \item Улучшает ли физически информированная Stage A/B/C калибровка прогнозную точность v3.5?
    \item Следует ли из лучшей прогнозной точности лучшая полезность суррогата как RL-среды?
\end{enumerate}

Ответы оказываются неодинаковыми. Суррогат v3 является компактной dual-head моделью с \(8\,482\) параметрами. Его задача --- не максимальная физическая точность на 24-часовом горизонте, а гладкая, дешёвая и устойчивая динамика для policy-gradient обучения. Суррогат v3.5 вводит физический backbone с параметром \(C_{\mathrm{zon}}\), который идентифицируется через Stage A/B/C inverse calibration. Внутри одной и той же 15-минутной выборки v3.5 после калибровки снижает 24-часовой RMSE с \SI{1.466}{\celsius} до \SI{0.644}{\celsius}. При этом matched-corpus эксперимент показывает, что большая часть исходного выигрыша объясняется переходом с hourly telemetry на 15-minute telemetry: v3 hourly даёт \SI{1.557}{\celsius}, v3 15-minute matched --- \SI{0.876}{\celsius}, а calibrated v3.5 --- \SI{0.644}{\celsius}.

Ключевой негативный результат возникает при переходе от offline prediction к live control. PPO, обученный напрямую на calibrated v3.5, проваливается на live BOPTEST: closed-loop RMSE составляет около \SI{4.32}{\celsius} на peak window и \SI{4.40}{\celsius} на typical window. Напротив, PPO на v3 остаётся существенно лучше, а hybrid backend даёт лучший компромисс: v3 обеспечивает rollout dynamics, а frozen v3.5 добавляет disagreement penalty. Это подтверждает центральный тезис: predictive fidelity является необходимым, но недостаточным критерием выбора surrogate backend для RL.

\begin{figure}[htbp]
    \centering
    \includegraphics[width=0.95\linewidth]{block1_q1_fig01_pipeline.pdf}
    \caption{Экспериментальный pipeline Блока 1: от BOPTEST telemetry и Stage A/B/C калибровки до live BOPTEST validation.}
    \label{fig:block1-pipeline}
\end{figure}

\section{Control-oriented суррогат v3}

\subsection{Назначение модели}

Суррогат v3 является рабочей моделью для обучения PPO. Его проектная цель --- не глобальная физическая точность, а быстрое и гладкое аппроксимирование локальной динамики, пригодное для миллионов rollout-шагов. В отличие от физически информированного v3.5, v3 не содержит явного \(C_{\mathrm{zon}}\), RC-backbone или ограничения на тепловой баланс. Это сознательный выбор: для policy-gradient алгоритма важны гладкие response surfaces, стабильные градиенты и отсутствие жёстких discontinuities.

Вход v3 является восьмимерным вектором признаков. Две головы используют один и тот же вход, но имеют независимые параметры:

\begin{itemize}
    \item \textbf{HeatFlowNetV2}: предсказывает приращение температуры зоны \(\Delta T\);
    \item \textbf{PowerNetV2}: предсказывает суммарную HVAC power \(P_{\mathrm{HVAC}}\).
\end{itemize}

\begin{table}[htbp]
\centering
\caption{Входной вектор v3: 8 признаков.}
\label{tab:v3-input}
\begin{tabular}{clll}
\toprule
Idx & Признак & Нормировка & Физический смысл \\
\midrule
0 & \(t_{\mathrm{zone,norm}}\) & \([15,35]^\circ C \rightarrow [-1,+1]\) & текущая температура зоны \\
1 & \(t_{\mathrm{amb,norm}}\) & \([-10,40]^\circ C \rightarrow [-1,+1]\) & наружная температура \\
2 & \(h_{\sin}\) & \(\sin(2\pi h/24)\) & суточная фаза \\
3 & \(h_{\cos}\) & \(\cos(2\pi h/24)\) & суточная фаза \\
4 & \(d_{\sin}\) & \(\sin(2\pi d/365)\) & сезонная фаза \\
5 & \(d_{\cos}\) & \(\cos(2\pi d/365)\) & сезонная фаза \\
6 & \(a_0\) & \([-1,+1]\rightarrow 18..35^\circ C\) & предыдущий \(T_{\mathrm{sup}}\) \\
7 & \(a_1\) & \([-1,+1]\rightarrow 0..1\) & предыдущая команда вентилятора \\
\bottomrule
\end{tabular}
\end{table}

\subsection{Архитектура v3}

Температурная голова использует блоки Linear+LayerNorm+Tanh и residual connection между первыми слоями. Power head является более простой MLP с Softplus-выходом, чтобы обеспечить \(P\geq 0\). Суммарный размер модели составляет \(8\,482\) обучаемых параметра: \(7\,105\) в температурной голове и \(1\,377\) в power head.

\begin{figure}[htbp]
    \centering
    \includegraphics[width=0.95\linewidth]{block1_q1_fig02_v3_dual_head.pdf}
    \caption{Dual-head архитектура v3: отдельные головы для \(\Delta T\) и HVAC power при общем восьмимерном входе.}
    \label{fig:v3-dual-head}
\end{figure}

\begin{table}[htbp]
\centering
\caption{Архитектура и training hyperparameters v3.}
\label{tab:v3-arch}
\begin{tabular}{ll}
\toprule
Свойство & Значение \\
\midrule
Heat head & Linear(8--64)+LN+Tanh \(\rightarrow\) Linear(64--64)+LN+Tanh+residual \(\rightarrow\) Linear(64--32)+Tanh \(\rightarrow\) Linear(32--1) \\
Power head & Linear(8--32)+Tanh \(\rightarrow\) Linear(32--32)+Tanh \(\rightarrow\) Linear(32--1)+Softplus \\
Activation & Tanh + LayerNorm + residual \\
Total params & \(8\,482\) \\
Optimizer & AdamW, lr \(=10^{-3}\), weight decay \(=10^{-4}\) \\
Scheduler & CosineAnnealingLR \\
Early stopping & patience 30, improvement threshold \(10^{-6}\) \\
Batch size & 256 \\
Multi-step horizons & [2,4] transition indices \\
\bottomrule
\end{tabular}
\end{table}

\subsection{Training trajectory и one-step validation}

Обучение запускалось максимум на 500 эпох, но остановилось на 215-й эпохе из-за early stopping. Лучший checkpoint соответствует эпохе 184. One-step validation RMSE составляет \SI{0.6255}{\celsius}, а \(R^2=0.9794\). Однако этот score не является достаточным доказательством control utility: split является contiguous и ограниченно репрезентативным, поэтому downstream claims опираются на external BOPTEST windows и rollout metrics.

\begin{table}[htbp]
\centering
\caption{Training result v3 checkpoint.}
\label{tab:v3-training}
\begin{tabular}{lll}
\toprule
Метрика & Значение & Источник \\
\midrule
Best epoch & 184 / 215 & checkpoint \\
Val RMSE one-step & \SI{0.6255}{\celsius} & checkpoint \\
Val \(R^2\) & 0.9794 & checkpoint \\
\bottomrule
\end{tabular}
\end{table}

\section{Predictive validity: multi-horizon rollout}

One-step RMSE не показывает, как модель ведёт себя в рекурсивном rollout. Поэтому основная метрика predictive validity --- multi-horizon rollout RMSE на горизонтах 1h, 4h, 8h и 24h.

\begin{table}[htbp]
\centering
\caption{Multi-horizon rollout RMSE: v3 против v3.5.}
\label{tab:multi-horizon}
\begin{tabular}{lccccc}
\toprule
Модель & 1h RMSE & 4h RMSE & 8h RMSE & 24h RMSE & 24h \(R^2\) \\
\midrule
v3 hourly & \SI{1.471}{\celsius} & \SI{1.579}{\celsius} & \SI{1.579}{\celsius} & \SI{1.557}{\celsius} & -1.411 \\
v3.5 raw & \SI{1.491}{\celsius} & -- & -- & \SI{1.466}{\celsius} & -- \\
v3.5 calibrated & \SI{0.321}{\celsius} & \SI{0.575}{\celsius} & \SI{0.638}{\celsius} & \SI{0.644}{\celsius} & -- \\
\bottomrule
\end{tabular}
\end{table}

\begin{figure}[htbp]
    \centering
    \includegraphics[width=0.95\linewidth]{block1_predictive_validity_horizon_lines.pdf}
    \caption{Multi-horizon rollout RMSE curves. Калиброванный v3.5 существенно снижает накопление ошибки на горизонтах 1--24 часа.}
    \label{fig:multi-horizon-rmse}
\end{figure}

Главный вывод из Table~\ref{tab:multi-horizon}: bare physical structure не гарантирует fidelity. Raw v3.5 лишь немного отличается от v3 на 24h horizon и может быть хуже matched v3 при контроле корпуса. Значительный выигрыш появляется только после Stage A/B/C calibration.

\section{Физически информированный суррогат v3.5 и Stage A/B/C}

\subsection{Мотивация v3.5}

v3.5 заменяет полностью black-box dynamics на физически структурированный backbone. Основной физический параметр --- тепловая ёмкость зоны:

\[
C_{\mathrm{zon}}.
\]

Идея состоит в том, что \(C_{\mathrm{zon}}\) является свойством здания, а не конкретной policy. Поэтому он должен быть интерпретируемым и потенциально переносимым между родственными BOPTEST testcases. Позже Block 3 подтверждает эту логику: hydronic-family testcases дают согласованные отношения \(C_{\mathrm{zon}}\) около \(1.9\times\) относительно \texttt{bestest\_air}.

\subsection{Stage A: preprocessing alignment}

Stage A очищает telemetry до физической идентификации. Он включает:

\begin{itemize}
    \item latency compensation между командой \(T_{\mathrm{sup}}\) и observed response;
    \item removal of constant temperature bias;
    \item affine normalization/scaling power channel.
\end{itemize}

\begin{table}[htbp]
\centering
\caption{Stage A preprocessing operations.}
\label{tab:stage-a}
\begin{tabular}{lll}
\toprule
Операция & Оценка & Эффект \\
\midrule
Latency compensation & 1 step = 15 min & RMSE search optimum \SI{0.3776}{\celsius} \\
Temperature bias removal & \(-\SI{0.1049}{\celsius}\) & median residual offset \\
Power affine normalization & least-squares fit & global scale + bias \\
Post-Stage-A alignment RMSE & \SI{0.2383}{\celsius} & baseline for Stage B \\
\bottomrule
\end{tabular}
\end{table}

Stage A является bookkeeping stage, а не learning stage. Его задача --- не позволить static bias или telemetry lag быть ошибочно интерпретированными как физический \(C_{\mathrm{zon}}\).

\subsection{Stage B: inverse identification of \(C_{\mathrm{zon}}\)}

Stage B оптимизирует \(C_{\mathrm{zon}}\) как скалярный физический параметр. Prior выбран как:

\[
C_{\mathrm{zon,prior}} = 4.200\times 10^5 \;\mathrm{J/K}.
\]

После 120 эпох Stage B получает:

\[
C_{\mathrm{zon,final}} = 4.413\times 10^5 \;\mathrm{J/K}.
\]

Относительное изменение:

\[
\frac{4.413-4.200}{4.200}\times 100\% \approx 5.1\%.
\]

\begin{table}[htbp]
\centering
\caption{Stage B identification of \(C_{\mathrm{zon}}\).}
\label{tab:stage-b}
\begin{tabular}{lll}
\toprule
Параметр & Значение & Комментарий \\
\midrule
\(C_{\mathrm{zon}}\) prior & \(4.200\times 10^5\) Дж/К & physical prior \\
\(C_{\mathrm{zon}}\) after Stage B & \(4.413\times 10^5\) Дж/К & \(+5.1\%\) vs prior \\
Stage B epochs & 120 & episode-aware run \\
Learning rate & 0.001 & scalar LR \\
Excitation quantile & 0.95 & top 5\% transient rows \\
Excitation rows & 403 / 8058 & high-identifiability subset \\
\bottomrule
\end{tabular}
\end{table}

\begin{figure}[htbp]
    \centering
    \includegraphics[width=0.95\linewidth]{block1_stage_abc_calibration_diagnostics.pdf}
    \caption{Stage B \(C_{\mathrm{zon}}\) convergence и Stage C residual-head diagnostics. График показывает, что физический параметр сходится, а не подгоняется вручную.}
    \label{fig:stage-abc-diagnostics}
\end{figure}

\subsection{Stage C: residual head refinement}

Stage C калибрует neural residual heads при frozen \(C_{\mathrm{zon}}\). В canonical pipeline Stage C реализован двухшагово:

\begin{enumerate}
    \item joint temperature/rollout/power refinement;
    \item power-head-only refinement без изменения \(C_{\mathrm{zon}}\) и temperature head.
\end{enumerate}

\begin{table}[htbp]
\centering
\caption{Stage C two-step calibration.}
\label{tab:stage-c}
\begin{tabular}{llll}
\toprule
Метрика & Step 1: episode-aware & Step 2: power-head-only & Combined effect \\
\midrule
Stage B epochs & 120 & 0 & \(C_{\mathrm{zon}}\) identified once \\
Stage C epochs & 60 & 80 & two-phase refinement \\
Temperature RMSE & \(0.3839\rightarrow 0.2348\) & unchanged & temperature alignment \\
Power MAE & \(810\rightarrow 808\) W & \(808\rightarrow 482\) W & power head refinement \\
\bottomrule
\end{tabular}
\end{table}

\subsection{Итоговая калибровочная сводка}

\begin{table}[htbp]
\centering
\caption{End-to-end v3.5 calibration summary.}
\label{tab:v35-summary}
\begin{tabular}{llll}
\toprule
Метрика & Raw / baseline & Calibrated & Improvement \\
\midrule
One-step RMSE\(_T\) & \SI{0.3839}{\celsius} & \SI{0.2348}{\celsius} & \(-38.9\%\) \\
24h rollout RMSE\(_T\) & \SI{1.4665}{\celsius} & \SI{0.6441}{\celsius} & \(-56.1\%\) \\
Power MAE & \SI{810.3}{W} & \SI{482.0}{W} & \(-40.5\%\) \\
\(C_{\mathrm{zon}}\) & \(4.200\times 10^5\) & \(4.413\times 10^5\) Дж/К & \(+5.1\%\) \\
\bottomrule
\end{tabular}
\end{table}

\begin{figure}[htbp]
    \centering
    \includegraphics[width=0.95\linewidth]{block1_q1_fig03_stage_abc_improvement.pdf}
    \caption{Effect of Stage A/B/C inverse calibration on v3.5 predictive fidelity.}
    \label{fig:stage-abc-improvement}
\end{figure}

\section{Matched-corpus control experiment}

Прямое сравнение legacy v3 и v3.5 было бы методологически слабым, потому что legacy v3 использовал hourly corpus, а v3.5 --- 15-minute corpus. Поэтому был выполнен matched-corpus control experiment: v3 был переобучен на том же 15-minute corpus, что и v3.5.

\begin{table}[htbp]
\centering
\caption{Matched-corpus comparison: 24h rollout RMSE.}
\label{tab:matched-corpus}
\begin{tabular}{lllll}
\toprule
Variant & Corpus & Step & Stage A/B/C & 24h RMSE\(_T\) \\
\midrule
v3 hourly & 51\,200 hourly transitions & 3600 s & n/a & \SI{1.5572}{\celsius} \\
v3 15-min matched & 10\,744 15-min transitions & 900 s & n/a & \SI{0.8761}{\celsius} \\
v3.5 raw & 10\,744 15-min transitions & 900 s & OFF & \SI{1.4665}{\celsius} \\
v3.5 calibrated & 10\,744 15-min transitions & 900 s & FULL & \SI{0.6441}{\celsius} \\
\bottomrule
\end{tabular}
\end{table}

Декомпозиция improvement:

\[
\Delta_{\mathrm{total}} = 1.5572 - 0.6441 = 0.9131^\circ C.
\]

Corpus shift contribution:

\[
\Delta_{\mathrm{corpus}} = 1.5572 - 0.8761 = 0.6811^\circ C \approx 74.6\%.
\]

Stage A/B/C contribution:

\[
\Delta_{\mathrm{calibration}} = 0.8761 - 0.6441 = 0.2320^\circ C \approx 25.4\%.
\]

\begin{figure}[htbp]
    \centering
    \includegraphics[width=0.95\linewidth]{block1_q1_fig06_fidelity_gain_waterfall.pdf}
    \caption{Attribution of v3-to-v3.5 predictive-fidelity gain: corpus shift vs Stage A/B/C calibration.}
    \label{fig:waterfall}
\end{figure}

Этот график является reviewer mitigation: статья не приписывает весь выигрыш физике. Улучшение является реальным, но bounded: основную часть даёт data-resolution shift, а оставшуюся часть --- физическая Stage A/B/C калибровка.

\section{Residual distribution и инженерные tolerance limits}

Средний RMSE не показывает форму ошибки. Поэтому дополнительно анализируются распределения residual:

\[
e_T = T_{\mathrm{pred}} - T_{\mathrm{obs}}.
\]

И абсолютная ошибка:

\[
|e_T| = |T_{\mathrm{pred}} - T_{\mathrm{obs}}|.
\]

\begin{figure}[htbp]
    \centering
    \includegraphics[width=0.95\linewidth]{block1_temperature_residual_histograms.pdf}
    \caption{Temperature residual distributions before and after calibration. Калибровка уменьшает spread и tail error.}
    \label{fig:residuals}
\end{figure}

Инженерная интерпретация residual plot:

\begin{itemize}
    \item смещение центра распределения показывает systematic bias;
    \item ширина распределения показывает типичную неопределённость;
    \item хвосты показывают редкие, но опасные ошибки;
    \item высокая доля samples в пределах \(\pm 1^\circ C\) означает пригодность для engineering tolerance.
\end{itemize}

\section{Runtime speed benchmark}

Скорость backend критична для RL, потому что PPO требует миллионы environment steps. Live BOPTEST RTE работает около 21 steps/s, тогда как v3 достигает около 4626 steps/s, v3.5 --- около 2399 steps/s, а hybrid backend --- около 1787 steps/s.

\begin{table}[htbp]
\centering
\caption{Backend speed benchmark.}
\label{tab:speed}
\begin{tabular}{lcc}
\toprule
Backend & Steps/s & Speed-up vs BOPTEST \\
\midrule
BOPTEST RTE & 21.0 & \(1\times\) \\
v3 & 4626.3 & \(220\times\) \\
v3.5 & 2399.5 & \(114\times\) \\
hybrid & 1786.8 & \(85\times\) \\
\bottomrule
\end{tabular}
\end{table}

\begin{figure}[htbp]
    \centering
    \includegraphics[width=0.85\linewidth]{block1_q1_fig09_backend_speed.pdf}
    \caption{Simulation throughput of BOPTEST and surrogate backends.}
    \label{fig:speed}
\end{figure}

Главный вывод: hybrid сохраняет достаточный speed-up для RL, несмотря на дополнительный frozen-v3.5 disagreement computation.

\section{Fidelity-to-RL gap: центральный негативный результат}

\subsection{Live BOPTEST architecture comparison}

Ключевой эксперимент Блока 1: те же surrogate variants используются как среды обучения PPO, а затем полученные policies проверяются в live BOPTEST RTE. Именно здесь возникает парадокс.

\begin{table}[htbp]
\centering
\caption{Live BOPTEST closed-loop transfer: surrogate fidelity vs RL utility.}
\label{tab:live-transfer}
\begin{tabular}{lcccc}
\toprule
Backend & Peak \(m_s\) & Typical \(m_s\) & Peak RMSE\(_T\) & Typical RMSE\(_T\) \\
\midrule
v3 & 0.0725 & 0.0947 & \SI{0.894}{\celsius} & \SI{0.745}{\celsius} \\
v3.5 calibrated direct & 1.0465 & 1.1020 & \SI{4.320}{\celsius} & \SI{4.401}{\celsius} \\
hybrid\_l010 & 0.0866 & 0.0411 & \SI{0.633}{\celsius} & \SI{0.612}{\celsius} \\
\bottomrule
\end{tabular}
\end{table}

\begin{figure}[htbp]
    \centering
    \includegraphics[width=0.90\linewidth]{block1_q1_fig07_fidelity_vs_rl_utility.pdf}
    \caption{Predictive fidelity does not imply RL training utility.}
    \label{fig:fidelity-vs-utility}
\end{figure}

Figure~\ref{fig:fidelity-vs-utility} является центральным графиком всей статьи. Он показывает, что calibrated v3.5 является лучшим predictive model, но худшим direct RL rollout backend. Следовательно, surrogate selection for RL cannot be based on predictive RMSE alone.

\subsection{Live controller performance}

\begin{figure}[htbp]
    \centering
    \includegraphics[width=0.92\linewidth]{block1_q1_fig08_live_boptest_performance.pdf}
    \caption{Live BOPTEST performance comparison of PPO controllers trained on different backends.}
    \label{fig:live-performance}
\end{figure}

Direct v3.5 PPO не просто немного хуже: он уходит в режим large comfort violation. Это означает, что ошибка не является statistical noise или minor tuning issue. Она связана с тем, что policy, обученная на v3.5 rollout dynamics, изучает action regimes, которые не переносятся в live BOPTEST.

\section{Failure mechanism: action saturation и transfer gap}

\subsection{Action saturation}

\begin{figure}[htbp]
    \centering
    \includegraphics[width=0.90\linewidth]{block1_q1_fig11_action_saturation.pdf}
    \caption{Policy action saturation under direct v3.5 training.}
    \label{fig:action-saturation}
\end{figure}

Direct v3.5 PPO демонстрирует bang-bang behavior: действия концентрируются около границ action space. Это объясняет, почему policy даёт плохую live performance, несмотря на хорошую offline predictive fidelity v3.5. Для RL важна не только точность прогнозирования, но и smoothness of action-response surface under the policy distribution.

\subsection{Transfer-gap diagnostics}

\begin{table}[htbp]
\centering
\caption{Transfer-gap diagnostics across backend variants.}
\label{tab:transfer-gap}
\begin{tabular}{lccc}
\toprule
Variant & \(m_s\) gap & Action gap norm & First divergence step \\
\midrule
pure v3 & \(-0.068\) to \(-0.102\) & 0.333--0.377 & 1 \\
direct v3.5 & \(-0.886\) to \(-1.014\) & 2.000--2.014 & 1 \\
hybrid\_l010 & \(-0.021\) to \(-0.023\) & 0.253--0.473 & 1--16 \\
\bottomrule
\end{tabular}
\end{table}

\begin{figure}[htbp]
    \centering
    \includegraphics[width=0.92\linewidth]{block1_q1_fig10_transfer_gap_diagnostics.pdf}
    \caption{Transfer-gap diagnostics reveal bang-bang saturation in standalone v3.5 training.}
    \label{fig:transfer-gap}
\end{figure}

Hybrid\_l010 резко уменьшает surrogate-to-live gap. Это объясняет, почему гибридная схема является конструктивным решением: она не использует v3.5 как direct environment, а применяет v3.5 как frozen physical censor.

\section{Hybrid backend: итоговое разрешение парадокса}

Гибридная функция награды имеет вид:

\[
r_{\mathrm{hybrid}} =
r_{\mathrm{comfort}} + r_{\mathrm{smooth}} + r_{\mathrm{energy}}
- \lambda_T |\Delta T_{\mathrm{disagree}}|
- \lambda_P |\Delta P_{\mathrm{disagree}}|.
\]

где:

\[
\Delta T_{\mathrm{disagree}} = T_{v3}(s,a) - T_{v3.5}(s,a),
\]

\[
\Delta P_{\mathrm{disagree}} = P_{v3}(s,a) - P_{v3.5}(s,a).
\]

\begin{figure}[htbp]
    \centering
    \includegraphics[width=0.95\linewidth]{block1_q1_fig14_hybrid_loss.pdf}
    \caption{Hybrid backend: v3 rollout dynamics with frozen-v3.5 soft regularization.}
    \label{fig:hybrid-loss}
\end{figure}

Гибридная схема разделяет роли:

\begin{itemize}
    \item v3 отвечает за rollout dynamics, smoothness и RL trainability;
    \item v3.5 отвечает за physical plausibility через disagreement penalty;
    \item PPO не оптимизируется напрямую в жёсткой v3.5 dynamics, которая вызывает action saturation.
\end{itemize}

\section{Replicative validity и residual diagnostics}

Для дополнительной проверки используются held-out BOPTEST episodes и residual distributions.

\begin{figure}[htbp]
    \centering
    \includegraphics[width=0.90\linewidth]{block1_q1_fig12_per_episode_rmse.pdf}
    \caption{Replicative validity across held-out BOPTEST episodes.}
    \label{fig:per-episode}
\end{figure}

\begin{figure}[htbp]
    \centering
    \includegraphics[width=0.90\linewidth]{block1_q1_fig13_residual_distributions.pdf}
    \caption{Temperature residual distributions before and after calibration.}
    \label{fig:residual-distributions}
\end{figure}

Эти дополнительные diagnostics подтверждают, что улучшение v3.5 не является следствием одной удачной trajectory или одного выбранного окна.

\section{Positioning относительно литературы}

\begin{figure}[htbp]
    \centering
    \includegraphics[width=0.95\linewidth]{block1_q1_fig15_literature_positioning.pdf}
    \caption{Positioning of this study against prior HVAC DRL and surrogate-model literature.}
    \label{fig:literature-positioning}
\end{figure}

Блок 1 занимает промежуточную позицию между literature on predictive digital twins and literature on RL control. Он показывает, что цифровой двойник должен оцениваться не только как predictor, но и как training environment. Именно это отличает текущую работу от исследований, где surrogate validation завершается на RMSE.

\section{Итоговые выводы Блока 1}

\subsection{Что доказано положительно}

\begin{enumerate}
    \item Stage A/B/C calibration действительно улучшает predictive fidelity v3.5.
    \item \(C_{\mathrm{zon}}\) идентифицируется физически правдоподобно: \(4.413\times 10^5\) Дж/К, \(+5.1\%\) к prior.
    \item Calibrated v3.5 снижает 24h rollout RMSE до \SI{0.644}{\celsius}.
    \item Hybrid backend сохраняет высокую скорость: около \(85\times\) быстрее BOPTEST RTE.
    \item Hybrid\_l010 даёт лучший live transfer compromise среди tested backends.
\end{enumerate}

\subsection{Что доказано отрицательно}

\begin{enumerate}
    \item Physical architecture without calibration is insufficient: raw v3.5 не является автоматически лучше v3.
    \item Better predictive fidelity does not imply better RL training utility.
    \item Direct PPO on calibrated v3.5 fails live BOPTEST despite strong offline RMSE.
    \item Action saturation является mechanism-level failure mode direct v3.5 training.
\end{enumerate}

\subsection{Центральный claim}

\begin{quote}
Более точный цифровой двойник не обязан быть лучшей средой обучения RL. В HVAC control необходимо разделять три роли: predictive fidelity, RL rollout utility и physical regularization. В нашем случае v3.5 наиболее полезен не как direct rollout environment, а как frozen physical teacher внутри hybrid backend.
\end{quote}

\subsection{Связь с последующими блоками}

Блок 1 даёт основу для Блока 2 и Блока 3:

\begin{itemize}
    \item Блок 2 проверяет, переносится ли hybrid role separation на разные controller families: thermostatic PPO, HDRL и MORL.
    \item Блок 3 проверяет, переносится ли Stage A/B/C surrogate recalibration на hydronic-family BOPTEST testcases.
\end{itemize}

Таким образом, Блок 1 не является isolated surrogate validation. Он формулирует главный принцип всей статьи: surrogate quality must be evaluated by function, not only by predictive score.

\end{document}
